\section{The polar triangle and side--angle duality}
\label{sec:polar}

The polar triangle is the mechanism behind the perfect duality of spherical
trigonometry.  In the present vector model it is constructive and explicit.

\begin{definition}[Positive polar triangle]
Define
\begin{equation}
 u^*=\frac{v\crossp w}{s_a},\qquad
 v^*=\frac{w\crossp u}{s_b},\qquad
 w^*=\frac{u\crossp v}{s_c}.                               \tag{5.1}
\end{equation}
The signs are fixed by
\(u\dotp u^*=v\dotp v^*=w\dotp w^*=\tau/(s_a,s_b,s_c)>0\),
so no binary choice of antipodes remains.
\end{definition}

Each polar vertex is a unit vector and is perpendicular to the great circle
containing the opposite original side.

\begin{theorem}[Polar side and angle correspondences]
\label{thm:polar-correspondence}
Let \(a^*,b^*,c^*\) and \(A^*,B^*,C^*\) denote the sides and angles of the
positive polar triangle.  Then
\begin{equation}
 a^*=\pi-A,
 \qquad b^*=\pi-B,
 \qquad c^*=\pi-C,                                         \tag{5.2}
\end{equation}
and
\begin{equation}
 A^*=\pi-a,
 \qquad B^*=\pi-b,
 \qquad C^*=\pi-c.                                         \tag{5.3}
\end{equation}
The polar operation is involutive: \((u^*)^*=u\), cyclically.
\end{theorem}

\begin{proof}
By Lagrange's four-vector identity,
\begin{align*}
 v^*\dotp w^*
 &=\frac{(w\crossp u)\dotp(u\crossp v)}{s_bs_c}\\
 &=\frac{(w\dotp u)(u\dotp v)-(w\dotp v)(u\dotp u)}{s_bs_c}\\
 &=\frac{yz-x}{s_bs_c}=-\cos A=\cos(\pi-A).
\end{align*}
The side \(a^*\) lies in \((0,\pi)\), and its sine is positive, so the branch is
uniquely \(a^*=\pi-A\).  The other side identities are cyclic.

For involution, use the vector triple-product identity:
\begin{equation}
 (w\crossp u)\crossp(u\crossp v)
 =u\bigl((w\crossp u)\dotp v\bigr)=\tau u.                 \tag{5.4}
\end{equation}
After division by \(s_bs_c\), the norm of the left-hand side is
\(\tau/(s_bs_c)=s_A=\sin a^*\).  Normalizing therefore returns \(u\).
The angle correspondence (5.3) follows by applying (5.2) to the polar triangle
and using involution.
\end{proof}

All steps in this proof are polynomial before positive normalization.  In
particular, polar duality does not require a compactness argument or a choice
of one of two normals: the orientation witness \(\tau>0\) fixes the sign.

\begin{corollary}[Duality principle]
\label{cor:duality-principle}
Any identity valid for all proper spherical triangles remains valid after the
simultaneous substitution
\begin{equation}
 (a,b,c;A,B,C)
 \longmapsto
 (\pi-A,\pi-B,\pi-C;\pi-a,\pi-b,\pi-c),                   \tag{5.5}
\end{equation}
with the corresponding changes
\(\sin(\pi-x)=\sin x\) and \(\cos(\pi-x)=-\cos x\).
\end{corollary}

Applying the side cosine law to the polar triangle immediately gives
\[
 -\cos A=\cos B\cos C-\sin B\sin C\cos a,
\]
which is Theorem~\ref{thm:angle-cosine}.  Thus the supplemental cosine law is
not an independent analytic theorem; it is the ordinary cosine law under the
finite polar construction.

\subsection{Half-sums under polarity}

Put
\begin{equation}
 p=\frac{a+b+c}{2},\qquad P=\frac{A+B+C}{2}.                \tag{5.6}
\end{equation}
For the polar triangle,
\begin{equation}
 p^*=\frac{a^*+b^*+c^*}{2}=\frac{3\pi}{2}-P,
 \qquad
 P^*=\frac{3\pi}{2}-p.                                    \tag{5.7}
\end{equation}
Consequently
\begin{align}
 \sin(p^*-b^*)&=\cos(P-B),&
 \sin p^*&=-\cos P,                                       \tag{5.8}\\
 \sin b^*&=\sin B,&
 \cos(P^*-B^*)&=\sin(p-b),                                \tag{5.9}
\end{align}
with cyclic companions.  These elementary substitutions turn the angle
half-formulas into the side half-formulas in Section~\ref{sec:half-formulas}.

\begin{corollary}[Polar excess and perimeter]
If \(E=A+B+C-\pi\) is the original spherical excess and \(E^*\) the polar
excess, then
\begin{equation}
 E^*=2\pi-(a+b+c).                                         \tag{5.10}
\end{equation}
Thus the proper-triangle condition \(a+b+c<2\pi\) is equivalent to positive
excess of the polar triangle.
\end{corollary}
